\documentclass[11pt,letterpaper]{article}
\usepackage[margin=1in]{geometry}
\usepackage{amsmath,amssymb}
\usepackage{xcolor}
\usepackage{booktabs}
\usepackage{mdframed}
\usepackage{enumitem}
\usepackage{fancyhdr}

\definecolor{ref}{HTML}{0C447C}
\definecolor{ans}{HTML}{0F6E56}
\definecolor{note}{HTML}{993C1D}

\newcommand{\refbox}[1]{\textcolor{ref}{\textit{[Suppl: #1]}}}
\newcommand{\ansbox}[1]{\colorbox{green!12}{\textbf{#1}}}
\newcommand{\ptag}[1]{\hfill\textit{(#1 pts)}}
\newcommand{\step}[1]{\medskip\noindent\textbf{Step #1.}\enspace}
\newcommand{\note}[1]{\par\noindent\textcolor{note}{\textit{Note: #1}}}

\pagestyle{fancy}\fancyhf{}
\renewcommand{\headrulewidth}{0.4pt}
\fancyhead[L]{\textbf{ECE332 --- Exam 1 Fall 2025 --- Answer Key}}
\fancyhead[R]{Bao Nguyen}
\fancyfoot[C]{\thepage}

\setlength{\parindent}{0pt}
\setlength{\parskip}{4pt}

\begin{document}

\begin{center}
{\LARGE\bfseries ECE332 Exam 1 --- Full Answer Key}\\[4pt]
{\large Fall 2025 $\cdot$ Step-by-Step with Supplement References}
\end{center}

\noindent\rule{\linewidth}{0.6pt}

%%%===========================================================
\section*{Problem 1 (20 pts) --- Transformer EMF, Triangular Loop}
%%%===========================================================

\textbf{Given:} Equilateral triangular loop, $N$ turns, side $a$, in $y$-$z$ plane with normal
$d\mathbf{s}=\hat{x}\,ds$. Magnetic field $\mathbf{B}(t)=B_0(-2\hat{x}+3\hat{y})\sin(\omega t)$.

\medskip

\textbf{(a) Flux through one turn.}\ptag{6}

\refbox{p.1 col 1: SURFACE ELEMENT $d\mathbf{s}$; DOT PRODUCTS; p.1 col 3: EMF \& FARADAY'S LAW}

\step{1} Find the area of the equilateral triangle.
From the exam cover: $A = \tfrac{1}{2}bh$.
For an equilateral triangle with side $a$: base $b=a$, height $h=\tfrac{\sqrt{3}}{2}a$.
\[
A = \tfrac{1}{2}\cdot a\cdot\tfrac{\sqrt{3}}{2}a = \tfrac{\sqrt{3}}{4}a^2
\]

\step{2} Compute flux. The loop is in the $y$-$z$ plane, so $d\mathbf{s}=\hat{x}\,ds$.
\[
\Phi = \iint_S \mathbf{B}\cdot d\mathbf{s}
= \iint B_0(-2\hat{x}+3\hat{y})\sin(\omega t)\cdot\hat{x}\,ds
\]
Use DOT PRODUCTS table: $\hat{x}\cdot\hat{x}=1$, $\hat{y}\cdot\hat{x}=0$.
\[
\Phi = B_0(-2)(1)\sin(\omega t)\iint ds + B_0(3)(0)\sin(\omega t)\iint ds
= -2B_0\sin(\omega t)\cdot A
\]
\[
\boxed{\Phi = -\frac{\sqrt{3}}{2}\,B_0\,a^2\sin(\omega t)\;\text{Wb}}
\]

\medskip

\textbf{(b) EMF induced in the loop.}\ptag{6}

\refbox{p.1 col 3: $V_\text{emf}=-N\,d\Phi/dt$; $\partial B/\partial t$ TABLE}

\step{1} Differentiate $\Phi$ with respect to $t$.
\[
\frac{d\Phi}{dt} = -\frac{\sqrt{3}}{2}\,B_0\,a^2\,\omega\cos(\omega t)
\quad\text{(from } \partial_t\sin(\omega t)=\omega\cos(\omega t)\text{ in }\partial B/\partial t\text{ table)}
\]

\step{2} Apply Faraday: $V_\text{emf}=-N\,d\Phi/dt$.
\[
V_\text{emf} = -N\left(-\frac{\sqrt{3}}{2}\,B_0\,a^2\,\omega\cos(\omega t)\right)
\]
\[
\boxed{V_\text{emf} = \frac{\sqrt{3}}{2}\,N\,B_0\,a^2\,\omega\cos(\omega t)\;\text{V}}
\]

\medskip

\textbf{(c) Current through 1\,k$\Omega$ resistor at given values.}\ptag{8}

\refbox{p.1 col 3: Equivalent circuit $I=V_\text{emf}/R$; p.1 col 2: UNIT CIRCLE $\cos(\pi/4)=\sqrt{2}/2$}

Plug in $N=10$, $a=0.08\,\text{m}$, $B_0=0.02\,\text{T}$, $\omega=2\times10^6\,\text{rad/s}$,
$\omega t=\pi/4$, $R=1\,\text{k}\Omega=10^3\,\Omega$.

\step{1} Combine $\cos(\pi/4)=\sqrt{2}/2$ with the $\sqrt{3}/2$ factor:
$\tfrac{\sqrt{3}}{2}\cdot\tfrac{\sqrt{2}}{2}=\tfrac{\sqrt{6}}{4}$. Then:
\[
V_\text{emf}
= \frac{\sqrt{6}}{4}\cdot N\cdot B_0\cdot a^2\cdot\omega
= \frac{\sqrt{6}}{4}\cdot 10\cdot 0.02\cdot(0.08)^2\cdot 2\times10^6
\]
Numeric bracket step by step:
\[
(0.08)^2=6.4\times10^{-3},\quad
0.02\times6.4\times10^{-3}=1.28\times10^{-4},\quad
1.28\times10^{-4}\times2\times10^6=256,\quad
10\times256=2560
\]
\[
V_\text{emf} = \tfrac{\sqrt{6}}{4}\times2560 = 640\sqrt{6}\approx 1568\;\text{V}
\]

\step{2} Since loop resistance is negligible ($R_i=0$), from EQUIVALENT CIRCUIT:
\[
I = \frac{V_\text{emf}}{R} = \frac{640\sqrt{6}}{10^3}
\]
\[
\boxed{I = 0.64\sqrt{6}\approx 1.57\;\text{A}}
\]

%%%===========================================================
\section*{Problem 2 (12 pts) --- Motional EMF, Moving Wire}
%%%===========================================================

\textbf{Given:} Wire of length $a=15\,\text{cm}=0.15\,\text{m}$ along $\hat{y}$.
Velocity $\mathbf{u}=12\hat{x}\,\text{cm/s}=0.12\hat{x}\,\text{m/s}$.
Field $\mathbf{B}=(0.1\hat{x}-0.05\hat{z})\,\text{T}$.

\refbox{p.1 col 3: Motional EMF $V^\text{m}=\oint(\mathbf{u}\times\mathbf{B})\cdot d\mathbf{l}$;
p.1 bottom: CROSS PROD CYCLE}

\step{1} Compute $\mathbf{u}\times\mathbf{B}$.
\[
\mathbf{u}\times\mathbf{B} = 0.12\hat{x}\times(0.1\hat{x}-0.05\hat{z})
= 0.012\underbrace{(\hat{x}\times\hat{x})}_{=\,0} - 0.006\underbrace{(\hat{x}\times\hat{z})}_{=\,-\hat{y}}
= 0 - 0.006(-\hat{y}) = 0.006\hat{y}\;\text{V/m}
\]

(From CROSS PROD CYCLE: $\hat{x}\times\hat{z}=-(\hat{z}\times\hat{x})=-\hat{y}$.)

\step{2} Integrate along the wire. Wire lies along $\hat{y}$, so $d\mathbf{l}=\hat{y}\,dy$.
\[
V_\text{emf} = \int_0^{0.15}(0.006\hat{y})\cdot(\hat{y}\,dy)
= 0.006\times 0.15
\]
\[
\boxed{V_\text{emf} = 9\times10^{-4}\;\text{V} = 0.9\;\text{mV}}
\]

\note{Only the $-0.05\hat{z}$ component of $\mathbf{B}$ contributes: $\hat{x}\times\hat{x}=0$
eliminates the $\hat{x}$ component entirely.}

%%%===========================================================
\section*{Problem 3 (12 pts) --- Conduction vs.\ Displacement Current}
%%%===========================================================

\textbf{Given:} $I_c=2\sin(\omega t)\,\mu\text{A}=2\times10^{-6}\sin(\omega t)\,\text{A}$ through area $A$.
Distilled water: $\sigma=10^{-4}\,\text{S/m}$, $\varepsilon_r=26\pi$,
$\varepsilon_0=\tfrac{1}{36\pi}\times10^{-9}\,\text{F/m}$.

\medskip

\textbf{(a) Electric field $E$.}\ptag{6}

\refbox{p.2 col 1: CONDUCTION vs.\ DISPLACEMENT CURRENT; $J_c=\sigma E$, $I_c=J_c A$}

\step{1} Current density: $J_c = I_c/A$.
\step{2} Ohm's law: $J_c=\sigma E$ $\Rightarrow$ $E=J_c/\sigma = I_c/(\sigma A)$.
\[
E = \frac{2\times10^{-6}\sin(\omega t)}{10^{-4}\cdot A}
\]
\[
\boxed{E = \frac{0.02}{A}\sin(\omega t)\;\text{V/m}}
\]

\medskip

\textbf{(b) Displacement current $I_d$ at $\omega=10^9\,\text{rad/s}$, $\omega t=\pi/3$.}\ptag{6}

\refbox{p.2 col 1: $J_d=\varepsilon\,\partial E/\partial t$, $I_d=J_d A$; UNIT CIRCLE: $\cos(\pi/3)=1/2$}

\step{1} Compute $\varepsilon=\varepsilon_r\varepsilon_0$.
\[
\varepsilon = 26\pi\cdot\frac{1}{36\pi}\times10^{-9} = \frac{26}{36}\times10^{-9}
= \frac{13}{18}\times10^{-9}\;\text{F/m}
\]

\step{2} $I_d = \varepsilon A\,\partial E/\partial t = \varepsilon A\cdot\frac{0.02}{A}\cdot\omega\cos(\omega t)
= \varepsilon\cdot 0.02\cdot\omega\cos(\omega t)$.

At $\omega t=\pi/3$:
\[
I_d = \frac{13}{18}\times10^{-9}\times 0.02\times10^9\times\frac{1}{2}
= \frac{13}{18}\times 0.02\times 0.5
= \frac{13}{1800}
\]
\[
\boxed{I_d = \frac{13}{1800}\approx 7.22\;\text{mA}}
\]

\note{$I_d\approx 7.22\,\text{mA}\gg I_c=2\,\mu\text{A}$. The water stores energy
capacitively ($\varepsilon_r=26\pi\approx82$) much more than it dissipates resistively at GHz.}

%%%===========================================================
\section*{Problem 4 (10 pts) --- Charge Relaxation in Distilled Water}
%%%===========================================================

\textbf{Given:} $\sigma\approx10^{-4}\,\text{S/m}$, $\varepsilon_r\approx26\pi$.

\medskip

\textbf{(a) Relaxation time constant.}\ptag{5}

\refbox{p.2 col 1: CHARGE RELAXATION $\tau_r=\varepsilon/\sigma$}

\[
\tau_r = \frac{\varepsilon_r\varepsilon_0}{\sigma}
= \frac{\tfrac{13}{18}\times10^{-9}}{10^{-4}}
= \frac{13}{18}\times10^{-5}\;\text{s}
\]
\[
\boxed{\tau_r = \frac{13}{18}\times10^{-5}\approx 7.22\;\mu\text{s}}
\]

\medskip

\textbf{(b) Time for $\rho_v$ to reach $e^{-3}$ of initial value.}\ptag{5}

\refbox{p.2 col 1: $\rho_v(t)=\rho_{v0}\,e^{-t/\tau_r}$}

\step{1} Set up equation:
\[
\rho_v(t)=\rho_{v0}\,e^{-t/\tau_r}=\rho_{v0}\,e^{-3}
\quad\Rightarrow\quad \frac{t}{\tau_r}=3 \quad\Rightarrow\quad t=3\tau_r
\]

\step{2} Compute:
\[
t = 3\times 7.22\;\mu\text{s}
\]
\[
\boxed{t = 3\tau_r\approx 21.7\;\mu\text{s}}
\]

%%%===========================================================
\section*{Problem 5 (30 pts) --- Lossless Plane Wave}
%%%===========================================================

\textbf{Given:} Nonconducting ($\sigma=0$) medium, $\varepsilon_r=6$, $\mu_r=1$.
\[
\mathbf{E}(z,t)=4\cos\!\left(6\pi\times10^8\,t - kz - \frac{\pi}{6}\right)\hat{y}\;\text{V/m}
\]

\medskip

\textbf{(a) Find $\omega$, $k$, $\lambda$, $\eta$.}\ptag{10}

\refbox{p.2 bottom: WAVE PARAMS; $k=\omega\sqrt{\mu\varepsilon}$, $\eta=\sqrt{\mu/\varepsilon}$, $\lambda=2\pi/k$}

\textbf{Read $\omega$ directly:}
\[
\boxed{\omega = 6\pi\times10^8\;\text{rad/s}}
\]

\textbf{Compute $k$:}
\[
k = \omega\sqrt{\mu_r\mu_0\,\varepsilon_r\varepsilon_0}
= \frac{\omega}{c}\sqrt{\mu_r\varepsilon_r}
= \frac{6\pi\times10^8}{3\times10^8}\sqrt{1\times6}
= 2\pi\sqrt{6}
\]
\[
\boxed{k = 2\pi\sqrt{6}\approx 15.4\;\text{rad/m}}
\]

\textbf{Compute $\lambda$:}
\[
\lambda = \frac{2\pi}{k} = \frac{2\pi}{2\pi\sqrt{6}} = \frac{1}{\sqrt{6}}
\]
\[
\boxed{\lambda = \frac{1}{\sqrt{6}}\approx 0.408\;\text{m}}
\]

\textbf{Compute $\eta$:}
\[
\eta = \sqrt{\frac{\mu}{\varepsilon}}
= \eta_0\sqrt{\frac{\mu_r}{\varepsilon_r}}
= 120\pi\sqrt{\frac{1}{6}}
= \frac{120\pi}{\sqrt{6}}
= 20\pi\sqrt{6}
\]
\[
\boxed{\eta = 20\pi\sqrt{6}\approx 154\;\Omega}
\]

\medskip

\textbf{(b) Find $\tilde{\mathbf{E}}(z)$.}\ptag{5}

\refbox{p.2 col 2: PHASOR METHOD step 2: $\cos(\omega t-kz+\phi)\to e^{j\phi}e^{-jkz}$}

Strip the $e^{j\omega t}$ factor. From $\mathbf{E}(z,t)=4\cos(\omega t-kz-\pi/6)\hat{y}$:
\[
\cos(\omega t - kz - \pi/6) \;\longleftrightarrow\; e^{-j\pi/6}e^{-jkz}
\]
\[
\boxed{\tilde{\mathbf{E}}(z) = 4\,e^{-j\pi/6}\,e^{-jkz}\,\hat{y}
= 4\,e^{-j(kz+\pi/6)}\,\hat{y}\;\text{V/m}}
\]
(with $k=2\pi\sqrt{6}$\,rad/m)

\medskip

\textbf{(c) Find $\tilde{\mathbf{H}}(z)$.}\ptag{5}

\refbox{p.2 col 2: DIRECTION RELATION $\tilde{\mathbf{H}}=\frac{1}{\eta}\hat{k}\times\tilde{\mathbf{E}}$;
table row: $\hat{k}=+\hat{z}$, $\hat{E}=+\hat{y}$ $\Rightarrow$ $\hat{H}=\hat{z}\times\hat{y}=-\hat{x}$}

\[
\tilde{\mathbf{H}} = \frac{1}{\eta}\hat{k}\times\tilde{\mathbf{E}}
= \frac{1}{\eta}\,\hat{z}\times\bigl(4e^{-j(kz+\pi/6)}\hat{y}\bigr)
= \frac{4}{\eta}\,e^{-j(kz+\pi/6)}\,(\hat{z}\times\hat{y})
= \frac{4}{\eta}\,e^{-j(kz+\pi/6)}\,(-\hat{x})
\]

With $\eta=20\pi\sqrt{6}$: $\tfrac{4}{\eta}=\tfrac{4}{20\pi\sqrt{6}}=\tfrac{1}{5\pi\sqrt{6}}$.
\[
\boxed{\tilde{\mathbf{H}}(z) = -\frac{1}{5\pi\sqrt{6}}\,e^{-j(kz+\pi/6)}\,\hat{x}\;\text{A/m}}
\]

\medskip

\textbf{(d) Find $\mathbf{H}(z,t)$.}\ptag{5}

\refbox{p.2 col 2: step 4: $\mathbf{H}=\operatorname{Re}\{\tilde{\mathbf{H}}\,e^{j\omega t}\}$}

\[
\mathbf{H}(z,t) = \operatorname{Re}\!\left\{-\frac{1}{5\pi\sqrt{6}}\,e^{-j(kz+\pi/6)}\,e^{j\omega t}\,\hat{x}\right\}
\]
\[
\boxed{\mathbf{H}(z,t)
= -\frac{1}{5\pi\sqrt{6}}\cos\!\left(6\pi\times10^8\,t - 2\pi\sqrt{6}\,z - \frac{\pi}{6}\right)\hat{x}\;\text{A/m}}
\]

\medskip

\textbf{(e) Average power density.}\ptag{5}

\refbox{p.2 col 3: POYNTING \& POWER (lossless): $\mathbf{S}_\text{av}=|E_0|^2/(2\eta)\,\hat{k}$}

\[
\mathbf{S}_\text{av} = \frac{|E_0|^2}{2\eta}\,\hat{k}
= \frac{4^2}{2\times 20\pi\sqrt{6}}\,\hat{z}
= \frac{16}{40\pi\sqrt{6}}\,\hat{z}
= \frac{2}{5\pi\sqrt{6}}\,\hat{z}
\]

Simplify: $\dfrac{2}{5\pi\sqrt{6}}\cdot\dfrac{\sqrt{6}}{\sqrt{6}}=\dfrac{2\sqrt{6}}{30\pi}=\dfrac{\sqrt{6}}{15\pi}$.
\[
\boxed{\mathbf{S}_\text{av} = \frac{\sqrt{6}}{15\pi}\,\hat{z}\approx 0.052\;\hat{z}\;\text{W/m}^2}
\]

%%%===========================================================
\section*{Problem 6 (16 pts) --- Lossy EM Wave in Distilled Water}
%%%===========================================================

\textbf{Given:} $f=\tfrac{2}{5}\,\text{GHz}=0.4\times10^9\,\text{Hz}$, $\sigma\approx10^{-4}\,\text{S/m}$,
$\varepsilon_r\approx26\pi$, $\mu_r=1$, wave in $+\hat{z}$, $|E(0)|=4\,\text{V/m}$.

Useful intermediate values (compute once, use everywhere):
\[
\omega = 2\pi f = 2\pi\times4\times10^8 = 8\pi\times10^8\;\text{rad/s}
\]
\[
\varepsilon' = \varepsilon_r\varepsilon_0
= 26\pi\cdot\frac{1}{36\pi}\times10^{-9}
= \frac{13}{18}\times10^{-9}\approx7.22\times10^{-10}\;\text{F/m}
\]

\medskip

\textbf{(a) Find $\varepsilon''/\varepsilon'$.}\ptag{6}

\refbox{p.2 col 3: MEDIUM TYPES; $\varepsilon''=\sigma/\omega$, $\varepsilon'=\varepsilon_r\varepsilon_0$}

\[
\varepsilon'' = \frac{\sigma}{\omega} = \frac{10^{-4}}{8\pi\times10^8}
= \frac{1}{8\pi}\times10^{-12}\;\text{F/m}
\]
\[
\frac{\varepsilon''}{\varepsilon'} = \frac{\sigma}{\omega\,\varepsilon_r\varepsilon_0}
= \frac{10^{-4}}{8\pi\times10^8\times\frac{13}{18}\times10^{-9}}
= \frac{10^{-4}}{\frac{104\pi}{18}\times10^{-1}}
= \frac{10^{-4}\times18}{104\pi\times10^{-1}}
= \frac{18}{104\pi}\times10^{-3}
= \frac{9}{52\pi}\times10^{-3}
\]
\[
\boxed{\frac{\varepsilon''}{\varepsilon'} = \frac{9}{52\pi}\times10^{-3}\approx 5.5\times10^{-5}}
\]

\medskip

\textbf{(b) What type of medium is distilled water at this frequency?}\ptag{5}

\refbox{p.2 col 3: MEDIUM TYPES table; threshold $\varepsilon''/\varepsilon'\ll 1$ (i.e.\ $<0.01$)}

\[
\frac{\varepsilon''}{\varepsilon'} \approx 5.5\times10^{-5} \;\ll\; 1
\]

\[
\boxed{\text{Low-loss dielectric}\quad\left(\frac{\varepsilon''}{\varepsilon'}\ll1\right)}
\]

Justification: the ratio is $\approx10^{-4}$, far below the $10^{-2}$ threshold for quasi-conductor
and far below 100 for good conductor.

\medskip

\textbf{(c) Find the skin depth.}\ptag{5}

\refbox{p.2 col 3: MEDIUM TYPES, low-loss row: $\alpha\approx\tfrac{\sigma}{2}\sqrt{\mu/\varepsilon'}$;
$\delta_s=1/\alpha$}

\step{1} Compute $\eta'=\sqrt{\mu_0/\varepsilon'}$ (intrinsic impedance).
\[
\eta' = \frac{\eta_0}{\sqrt{\varepsilon_r}}
= \frac{120\pi}{\sqrt{26\pi}}
\approx \frac{376.99}{9.04}
\approx 41.7\;\Omega
\]

\step{2} Apply low-loss approximation $\alpha\approx(\sigma/2)\,\eta'$:
\[
\alpha \approx \frac{\sigma}{2}\cdot\eta'
= \frac{10^{-4}}{2}\times41.7
= 2.085\times10^{-3}\;\text{Np/m}
\]

\step{3} Skin depth $\delta_s=1/\alpha$:
\[
\delta_s = \frac{1}{2.085\times10^{-3}}
\]
\[
\boxed{\delta_s \approx 480\;\text{m}}
\]

\note{480\,m skin depth confirms this is a near-lossless medium at $0.4$\,GHz with $\sigma=10^{-4}$\,S/m.
The wave barely attenuates --- it can travel hundreds of meters before losing a factor of $e$.}

\bigskip
\noindent\rule{\linewidth}{0.6pt}
\begin{center}
\textit{Blue labels \refbox{...} point to the exact section of the Exam 1 Supplement to look up.}
\end{center}

\end{document}
